\newcommand{\CourseTemplatePath}{../../../../../templates/slides/}
\input{\CourseTemplatePath course_preamble.tex}
\title[Course Introduction]{Course Introduction}
\subtitle{Aims, syllabus, learning activities, and assessment}
\begin{document}
\makecoursetitle

% Suggested presentation time: 2 minutes.
\begin{frame}{What Will You Learn?}
Use computation to turn an economic model into results that you can explain and check.
\medskip
\begin{itemize}
\item Translate a model into states, decisions, equations, and an algorithm.
\item Choose numerical methods and understand their assumptions and approximations.
\item Implement a calculation, check its output, and make it reproducible.
\item Read computational research and connect its methods to your own questions.
\end{itemize}
\end{frame}

% Suggested presentation time: 1.5 minutes.
\begin{frame}{How the Module Works}
\begin{tabularx}{\linewidth}{@{}lY@{}}
\toprule\textbf{Activity} & \textbf{What you will do}\\\midrule
Lectures & Build intuition, define the mathematical objects, and study computational methods\\
Labs & Work with practical tools, run examples, and verify results\\
Homework & Read, practice, and explain what you have learned\\
Final project & Apply computational methods through a research-paper replication\\\bottomrule
\end{tabularx}
\par\medskip
Each class has two parts of about 50 minutes. Weekly pages organize the schedule; topic-based decks may continue across sessions.
\end{frame}

% Suggested presentation time: 2 minutes.
\begin{frame}{Syllabus at a Glance}
\small
\begin{tabularx}{\linewidth}{@{}lY@{}}
\toprule\textbf{Sessions} & \textbf{Topics}\\\midrule
Opening & Course introduction; Markov processes and Markov decision processes\\
Core methods & Dynamic programming; Euler equations and endogenous grids; expectations; simulation\\
Approximation and learning & Function approximation; neural networks; supervised and unsupervised learning; reinforcement learning\\
Applications and synthesis & Heterogeneous-agent models; connections between methods\\
Research literature & Computational economics using machine-learning methods\\\bottomrule
\end{tabularx}
\par\medskip
See the \href{https://yaolangzhong.github.io/U_Tokyo_Comp_Econ_Course/syllabus.html}{online syllabus} for the weekly schedule. Literature selections and the closing-session programme will be announced there.
\end{frame}

% Suggested presentation time: 2 minutes. Let students scan the code.
\begin{frame}{Your Course Website}
\begin{columns}[T,onlytextwidth]
\begin{column}{0.70\textwidth}
\href{https://yaolangzhong.github.io/U_Tokyo_Comp_Econ_Course/}{\textbf{Open Computational Economics}}
\begin{itemize}
\item \textbf{Syllabus:} class dates and topic sequence.
\item \textbf{Weekly materials:} lecture and lab slides, notebooks, and homework decks.
\item \textbf{References:} reading resources and source links.
\item \textbf{Grading:} assessment and replication-project requirements.
\end{itemize}
\small Bookmark the homepage. Use \textbf{Open in Colab} to run the introductory notebook online, or download it for VS Code.
\end{column}
\begin{column}{0.25\textwidth}
\centering
\href{https://yaolangzhong.github.io/U_Tokyo_Comp_Econ_Course/}{\includegraphics[width=\linewidth]{../../../../../assets/course-qr.png}}
\par\smallskip\footnotesize Scan to open the course website.
\end{column}
\end{columns}
\end{frame}

% Suggested presentation time: 1.5 minutes. Grading policy awaits confirmation.
\begin{frame}{Assessment and Grading}
The module includes a replication project and weekly reading or exercise activities.
\par\medskip
\begin{tabularx}{\linewidth}{@{}lY@{}}
\toprule\textbf{Assessment detail} & \textbf{Announcement}\\\midrule
Grade weights & To be announced\\
Homework assessment & To be announced\\
Project team size and scope & To be announced\\
Deadlines and submission channel & To be announced\\\bottomrule
\end{tabularx}
\par\medskip
The confirmed assessment arrangements will appear in the syllabus and project guide.
\end{frame}

% Suggested presentation time: 1.5 minutes.
\begin{frame}{The Replication Project}
Choose a research paper and reproduce its computational or quantitative methods.
\begin{itemize}
\item Explain the research question, key equations, and algorithm.
\item Organize runnable code with clear setup and execution instructions.
\item Check your results and explain differences from the paper.
\item Present your replication work in a \textbf{15-minute presentation}.
\end{itemize}
\medskip
Presentations take place in the last one or two lecture weeks, depending on the number of submissions. See the \href{https://yaolangzhong.github.io/U_Tokyo_Comp_Econ_Course/project.html}{grading and project guide}.
\end{frame}

% Suggested presentation time: 1.5 minutes.
\begin{frame}{Our Opening Session}
\begin{tabularx}{\linewidth}{@{}lY@{}}
\toprule\textbf{Time} & \textbf{Activity}\\\midrule
First part: about 15 minutes & Course introduction: aims, website, syllabus, and assessment\\
First part: about 35 minutes & Introduction to Markov decision processes (MDP), beginning with Markov processes\\
Second part: about 50 minutes & Introduction to Git and Python; development tools and a first notebook\\\bottomrule
\end{tabularx}
\par\medskip
After class, use the separate \href{https://yaolangzhong.github.io/U_Tokyo_Comp_Econ_Course/downloads/week_01/homework/Setting_Up_Git_and_Python.pdf}{Setting Up Git and Python} homework deck for the setup exercise.
\end{frame}

% Suggested presentation time: 1.5 minutes, including questions.
\begin{frame}{Practical Information and Questions}
\begin{itemize}
\item \textbf{Teaching period:} October 7, 2026--January 27, 2027. Consult the online schedule for individual class dates.
\item \textbf{Meeting:} Wednesdays, Period 4; International Academic Research Building, Seminar Room 514.
\item \textbf{Contact:} Yaolang Zhong, \href{mailto:yaolang.zhong@e.u-tokyo.ac.jp}{yaolang.zhong@e.u-tokyo.ac.jp}.
\item \textbf{Office hours and consultation arrangements:} to be announced.
\end{itemize}
\medskip
Questions about the course or preparation?\par
Next: examples and definitions of Markov processes.
\end{frame}
\end{document}
